\clearpage
\subsection{Picard--Lefschetz 数据}\label{sec:19c}
再次考虑 $(n+1)$ 维光滑射影簇，这次明确将它嵌入射影空间 $\overline X\subset\CP^N$。与前面一样，对某个 $d\in\Z$，要求
\begin{equation}\label{eq:19.5}
\mathcal K_{\overline X}\cong\mathcal O_{\overline X}(-d).
\end{equation}
假设射影嵌入相对于标准齐次坐标 $(y_0:\cdots:y_N)$ 处于一般位置。准确的条件将在下面列出，但要点是，施加一个适当的 $\operatorname{PGL}_{N+1}$ 元素后，或者等价地，换用另一齐次坐标系后，这些条件总能成立。

取无穷远超平面为 $\CP^{N-1}=\{y_N=0\}$，将其补集等同于 $\CC^N$，坐标为 $(x_0=y_0/y_N,\ldots,x_{N-1}=y_{N-1}/y_N)$。对 $0\le k\le n+1$，取子空间 $H^k=\{x_0=\cdots=x_{k-1}=0\}\subset\CC^N$ 及其闭包 $\overline H^k\subset\CP^N$。定义
\[
\begin{aligned}
\overline X^k&=\overline X\cap\overline H^k, &X^k&=X\cap H^k,\\
\Psi^k&=x_k:X^k\longrightarrow\CC,&
\varpi^k&=(x_k,x_{k+1}):X^k\longrightarrow\CC^2.
\end{aligned}
\]
一般情形下，以下性质成立：
\begin{itemize}
\item 对 $0\le k\le n+1$，$\overline X$ 与全部 $\overline H^k$ 及 $\overline H^k\cap\CP^{N-1}$ 横截。这意味着每个 $X^k$ 光滑，而且在无穷远处也光滑。
\end{itemize}
注意，按定义，$X^{k+1}=(\Psi^k)^{-1}(0)$、$\Psi^{k+1}=\varpi^k|X^{k+1}$。此外，由\eref{19.5}，有 $\mathcal K_{\overline X^k}\cong\mathcal O_{\overline X^k}(k-d)$。\znote{299}{这是本节最容易错位的指标约定：$\Psi^k=x_k$ 的纤维是 $X^{k+1}$，所以由 $\Psi^k$ 的临界点集合 $C^k$（含于 $X^k$）标号的消失圈都住在 $X^{k+1}$ 里，$\A^{k+1}$ 才是由它们生成的有向范畴。\eref{19.6} 中 $\F(X^k)\to\F(\Psi^k)\to\Tw\A^{k+1,\to}$ 的指标比 $k$ 大一位，正是由此而来。原书 p302 有一处把 $\gamma_j^{k-1}$ 印作 $\gamma_j^{k+1}$，读时须按此约定校正。顺带一提，$\varpi^k=(x_k,x_{k+1})$ 的两个分量正是 $\Psi^k$ 与 $\Psi^{k+1}$，故双纤维化的指标同样是“错一位”的。[来源：A，原书 (19b)–(19c) 的定义与 (19.6)。]}继续列出一般性质：
\begin{itemize}
\item 对 $0\le k\le n$，每个函数 $\Psi^k$ 只有非退化临界点，且任意两点不在同一纤维中。
\SourcePage{308}
\item 对 $0\le k\le n-1$，$D\varpi^k$ 是一般的，即横截于 $\Hom(TX^k,\CC^2)$ 中由低秩线性映射构成的各子簇。若 $w\in\CC$ 为 $\Psi^k$ 的临界值，$x$ 为纤维 $(\Psi^k)^{-1}(w)$ 中的唯一临界点，则还满足：$(D^2\Psi^k)_x$ 在 $(D\varpi^k)_x$ 的核上的限制仍非退化；函数
\[
(\Psi^k)^{-1}(w)\setminus\{x\}\xrightarrow{\ \varpi^k\ }\{w\}\times\CC=\CC
\]
具有非退化临界点；任意两个这样的点的值不同，并且也都不同于 $\varpi^k(x)$ 本身。
\end{itemize}
这些性质在一般情形下成立是经典事实，但为方便起见，仍概述证明。第一部分是 Bertini--Sard 定理的直接应用。第二部分可用以下标准论证归约到同一定理。考虑拉回丛 $T^*X^k\to\CC^{N-k}\times X^k$ 的截面 $(\lambda,x)\mapsto d(\sum_{i\ge k}\lambda_i x_i)$。它总是与零截面横截相交，所以其零点集 $Z\subset\CC^{N-k}\times X^k$ 光滑。若取投影 $Z\to\CC^{N-k}$ 的一个正则值 $\lambda$，则函数 $\sum_i\lambda_i x_i$ 只有非退化临界点。显然，任意这样的函数都可通过线性坐标变换变成 $x_k$。因此，给定任意 $\overline X$，可用一个接近恒等的 $\operatorname{PGL}_{N+1}$ 元素改变其嵌入，使变换后的簇上函数 $\Psi^k=x_k$ 只有非退化临界点。再对嵌入作微扰，可使各临界点位于不同纤维。同一论证适用于 $\varpi^k$。读者想必已经注意到，这些条件是\secref{15e}中条件的代数几何对应物。其推论也类似：若 $x$ 为 $\Psi^k$ 的临界点，值为 $w$，则 $\operatorname{Critv}(\varpi^k)$ 在 $\varpi^k(x)$ 处与直线 $\{w\}\times\CC$ 相切，在其他各处都与之横截。按\secref{15e}，定义 $C^k=\operatorname{Critv}(\Psi^k)$，$0\le k\le n$，及 $F^k=\operatorname{Fakev}(\Psi^k,\varpi^k)$，$0\le k\le n-1$。为简化记号，令 $F^n=\varnothing$。注意，由假设，$0\notin C^k\cup F^k$。

固定 $0\le k\le n$，写 $C^k=\{c_j^k\}$。为每个 $j$ 选择嵌入路径 $\gamma_j^k:[0,1]\to\CC$，满足 $\gamma_j^k(0)=0$、$\gamma_j^k(1)=c_j^k$，$t<1$ 时 $\gamma_j^k(t)\notin C^k$，且对所有 $t$，$\gamma_j^k(t)\notin F^k$。这些路径应构成\secref{16d}意义下的基。在最低维情形 $k=n$，它们只是从 $0$ 走向分歧覆盖 $\Psi^n:X^n\to\CC$ 的某个分歧点的路径。可为每条路径配上无序点对 $\mu_j^{n+1}\subset X^{n+1}$，标记沿路径合并的覆盖的两层。更高维时，如\secref{16i}那样，对 $t\in[0,1)$，追踪子集 $\operatorname{Critv}(\varpi^k)\cap(\{\gamma_j^k(t)\}\times\CC)\subset\CC$ 的移动。由此得到嵌入路径 $\mu_j^{k+1}:[-1,1]\to\CC$，满足 $(\mu_j^{k+1})^{-1}(C^{k+1})=\{-1,1\}$，在同痕及反向意义下唯一。这一过程产生以下信息，称为 $X$ 的 \emph{Picard--Lefschetz 数据}：
\SourcePage{309}
\begin{itemize}
\item 有限集 $X^{n+1}$，以及子集族 $\{\mu_j^{n+1}\}$；
\item 对每个 $0\le k\le n$，在带标记点 $\{0\}\cup C^k\cup F^k$ 的平面 $\CC$ 中的路径集 $\{\gamma_j^k\}$，以及 $k>0$ 时的 $\{\mu_j^k\}$；
\item 当 $H^1(\overline X)\ne0$ 时，还需一个稍后说明的同伦数据。
\end{itemize}
注意，这些应被视为拓扑图形：可移动标记点，同时对相关路径作适当同痕，而不实质改变数据。因此，可通过指定一个穿孔平面微分同胚如何把标准路径变成给定路径，来编码每条路径。这样，剩下的主要是一组辫群中的字，这是 Moishezon 辫分解技巧的弱化版本。

考察这些数据包含多少关于 $X$ 的信息。首先，分歧数据容许在拓扑上重构 $\Psi^n:X^n\to\CC$。为了得到 $X^n$ 的良好定义的 Fukaya 范畴，还需知道 $\mathcal K_{X^n}$ 的平凡化，至少在拓扑上知道，即知道其在同伦意义下的 $C^\infty$ 平凡化。由构造，该平凡化来自一个处处非零的复一形式，在每个无穷远点处有 $d-n$ 阶极点。为确定其同伦类，应记录一组张成 $H_1(\overline X^n)$ 的浸入曲线的“旋转数”，即 Maslov 指标。若一条曲线是 $\Psi^{n-1}$ 的消失圈，已知其旋转数为零，因为该曲线容许分次。假设 $H^1(\overline X)=0$。由弱 Lefschetz 定理，$H^1(\overline X^{n-1})=0$，且消失圈张成 $H_1(\overline X^n)$，这意味着 $\mathcal K_{X^n}$ 的平凡化唯一确定。否则，可能还须记录一些额外旋转数，这就是上面提到的最后一项数据。

\StatementNumber{9}\begin{note}\label{st:19.9}
另一种做法有些偏离一般思路：假设精确知道分歧点集 $C^n$。此时，$X^n$ 在全纯同构意义下确定，并带有一个在非零常数倍意义下唯一的全纯一形式，它没有零点，在无穷远点附近具有所需阶数的极点。
\end{note}

在上一层，$X^n$ 中存在唯一的简单闭曲线，使它们在 $\Psi^n$ 下的像为 $\mu_j^n$。此外，这些曲线都不与零同调，因为它们来自消失圈，故对适当的一形式选择是正合的。按 Donaldson 与 Gompf 的方法，这组曲线连同相应路径 $\gamma_j^{n-1}$，在辛同构意义下确定 $\overline X^{n-1}$；这里辛同构相对于无穷远除子而言，因而诱导 $X^{n-1}$ 的正合辛同构。最后，$\mathcal K_{X^n}$ 的平凡化以在同伦意义下唯一的方式，诱导 $\mathcal K_{X^{n-1}}$ 的平凡化。

对 $k<n$，Picard--Lefschetz 数据本身是否足以确定 $X^k$ 或 $\overline X^k$ 的辛同构类型，似乎值得怀疑。因为我们记录了匹配路径，却没有记录构造实际
\SourcePage{310}
匹配圈所需的同痕。将 $X^k$ 紧化为 $\overline X^k$ 时，实际上还有一个潜在问题，即无穷远除子的附着方式。可尝试通过指定额外数据克服这些问题，但我们不采用这一方向。\znote{302}{为什么说数据不足，可与上文三级信息对照：$X^n$ 由分歧数据与旋转数确定；$X^{n-1}$ 由 $X^n$ 中的曲线连同路径 $\gamma_j^{n-1}$ 确定；再往下一层，则需把 $\gamma_j^{k-1}$ 同痕成匹配路径 $\mu_j^k$ 的具体同痕，而同痕不同可以使得到的匹配圈相差一次 Lagrangian 手术，$X^k$ 的辛结构随之改变。紧化还多一层歧义：无穷远除子如何附着并未被数据记录。[来源：C，译者按（上下文推断，无外部依据），据本节末三段。]}
